\chapter{Considerações finais}

O presente estudo buscou explorar a relação entre \textit{Branding}, sustentabilidade e empreendedorismo sustentável de empregos agroflorestais na região amazônica, abordando especificamente a percepção de potenciais consumidores e os empreendedores no contexto de marcas associadas à proteção ambiental e desenvolvimento sustentável. Através do desenvolvimento e validação de uma escala psicometricamente sólida, este trabalho contribuiu significativamente para a compreensão e aplicação das práticas de \textit{Branding} sustentável no cenário empresarial amazônico.

Os resultados deste estudo confirmam a importância do \textit{Branding} sustentável como uma ferramenta de diferenciação e valorização de produtos e marcas relacionadas à conservação florestal. A percepção positiva dos consumidores frente a produtos sustentáveis, juntamente com a disposição para pagar um preço premium, indica um mercado crescente e promissor para iniciativas de empreendedorismo sustentável na Amazônia.

A escala BES Amazônia proposta, revelou-se um instrumento valioso para a análise da percepção dos consumidores, bem como para a avaliação da gestão de \textit{Branding} de marcas sustentáveis. Sua aplicação em diferentes grupos populacionais e em contexto de pesquisa de mercado oferece aos gestores uma perspectiva clara e objetiva das preferências e expectativas dos consumidores em relação às marcas sustentáveis, permitindo o refinamento das estratégias de marketing e comunicação.

A inclusão da autenticidade e da conexão entre consumidores e marcas em estudos futuros deverá ampliar ainda mais nossa compreensão sobre as nuances envolvidas no \textit{Branding} sustentável. Esses elementos, juntamente com a lealdade à marca e o compromisso dos consumidores, são fundamentais para criar e fortalecer relacionamentos de longo prazo com o público, garantindo a continuidade e o sucesso dos empreendimentos sustentáveis na Amazônia.

Este estudo também destaca a importância do engajamento e colaboração entre todos os stakeholders envolvidos no processo de empreendedorismo sustentável, incluindo a interação entre empresas, comunidades tradicionais, agentes governamentais e consumidores. O desenvolvimento de estratégias participativas e cooperativas, que busquem harmonizar interesses socioeconômicos e ambientais, é crucial para superar os desafios e garantir a prosperidade a longo prazo dos negócios sustentáveis na região amazônica.

Este estudo é uma base para a pesquisa em \textit{Branding} sustentável e empreendedorismo na região amazônica, que pode servir como um ponto de partida para futuros estudos nessa área. Dada a importância da Amazônia no equilíbrio ambiental global e a crescente conscientização dos consumidores em relação às questões de sustentabilidade, é fundamental que os gestores de marcas compreendam os fatores que moldam a percepção, atitudes e comportamento do público em relação a tais produtos e serviços.

Além disso, este estudo destaca a necessidade de desenvolver estratégias de marketing mais direcionadas, que enfatizem os atributos sustentáveis dos produtos e serviços oferecidos, bem como a importância de envolver e colaborar com stakeholders locais, para garantir a continuidade e a legitimidade de tais iniciativas no contexto amazônico.
Futuras pesquisas poderiam explorar a relação entre \textit{Branding} sustentável e empreendedorismo em outros contextos geográficos e culturais, buscando identificar semelhanças e diferenças nas percepções e comportamentos dos consumidores. Além disso, estudos longitudinais e comparativos poderiam avaliar a eficácia de diferentes estratégias de marketing e comunicação na promoção da sustentabilidade e adoção de práticas sustentáveis pelos consumidores ao longo do tempo.

Outra linha de pesquisa potencial pode considerar a avaliação da efetividade de diferentes formatos e abordagens de comunicação no que diz respeito ao engajamento dos consumidores e à construção de relações sólidas entre marca e público. Além disso, o papel dos líderes de opinião, influenciadores digitais e grupos de interesse poderia ser investigado em maior profundidade, para melhor compreender seus impactos e influências na percepção e comportamento dos consumidores em relação às marcas e produtos sustentáveis.

Também é importante reconhecer o papel da educação e conscientização ambiental na formação de opiniões e atitudes em relação à sustentabilidade e \textit{Branding} sustentável. Portanto, estudos adicionais poderiam examinar a influência de campanhas educacionais e de conscientização na construção de uma representação mental positiva da sustentabilidade e na disposição dos consumidores para apoiar e adotar práticas sustentáveis em suas vidas cotidianas.

As principais dificuldades percebidas pelos empreendedores locais na negociação de produtos sustentáveis no mercado, evidenciadas neste estudo, indicam a necessidade de uma maior atenção às estratégicas de comunicação e marketing, de modo a destacar as vantagens e benefícios das práticas sustentáveis e ambientalmente responsáveis. A adoção de signos e certificações que comprovem a sustentabilidade das marcas é uma estratégia-chave para promover a confiança dos consumidores e impulsionar as atividades econômicas de tais empreendimentos.

Enfatiza-se, a relevância deste estudo como base para futuras pesquisas na área de \textit{Branding} sustentável e empreendedorismo na Amazônia. Podendo ser base para novos ensaios tais como: análises comparativas entre diferentes culturas e países, investigando as percepções e atitudes dos consumidores sobre marcas sustentáveis oriundas da Amazônia do Brasil em diferentes contextos culturais e internacionais, permitindo uma compreensão mais ampla das variáveis que influenciam a percepção e comportamento dos consumidores em relação a tais marcas.

Estudos de caso em empreendedorismo sustentável que examine outras iniciativas bem-sucedidas e inovadoras de empreendedorismo sustentável na região amazônica e outras cadeias produtivas sustentáveis, a fim de identificar outras práticas, estratégias e modelos de negócios que possam ser replicadas ou adaptadas suas estratégias de \textit{Branding} para promover o desenvolvimento sustentável em outros contextos.

O impacto da tecnologia e inovação no \textit{Branding} sustentável, investigando como os avanços tecnológicos, tais como a inteligência artificial e a Internet das Coisas (IoT), podem influenciar e aprimorar estratégias de \textit{Branding} e marketing sustentáveis, aumentando a eficiência dessas práticas e facilitando a coleta e análise de dados relacionados ao engajamento dos consumidores;
A influência das redes sociais, o marketing e os influenciadores digitais influenciam a promoção de marcas e produtos sustentáveis, especialmente nas gerações mais jovens, que são mais propensas a utilizar essas plataformas e seguir tais influenciadores.

A Avaliação do desempenho e impacto social dos empreendimentos sustentáveis, abordando o desempenho financeiro e o impacto social de empresas sustentáveis na região amazônica, analisando também os desafios e benefícios específicos para esses empreendedores e comparando-os com empresas em outros contextos geográficos e setores da economia.

E por fim, as políticas públicas e governança no empreendedorismo sustentável, examinando o papel das políticas públicas, regulações e incentivos governamentais no desenvolvimento de empreendimentos sustentáveis na Amazônia, assim como identificar melhores práticas e oportunidades de cooperação entre os setores público, privado e não governamental em prol da sustentabilidade.
Ao abordar tais tópicos em pesquisas futuras, acadêmicos e profissionais poderão obter percepções valiosos sobre as melhores práticas e oportunidades de desenvolvimento do \textit{Branding} sustentável e do empreendedorismo na região amazônica e em outros contextos.
 

Com o aprofundamento do conhecimento nessa área, é possível serem identificadas maneiras mais eficazes de promover a sustentabilidade e preservação ambiental ao nível global, ao mesmo tempo, em que se apoia no desenvolvimento econômico equitativo e a prosperidade das comunidades afetadas.
